% Instructor-authored non-Hermitian physics interlude (2026-10-17).
% Complete worked model immediately after finite-dimensional Hermiticity.
\section{Beyond Hermiticity: Loss, Complex Spectra, and Exceptional Points}
\label{sec:nh-two-cavities}
A Hermitian Hamiltonian has real eigenvalues and preserves the norm of a
state under Schr\"odinger evolution. A subsystem that exchanges
excitations with its environment admits another useful description:
an \emph{effective non-Hermitian Hamiltonian}. We now construct one
directly from a physical model, write down every matrix element, and
follow its eigenvalues and dynamics without assuming diagonalizability.

\subsection*{1) What changes when a Hamiltonian is non-Hermitian?}
Let the possibly non-Hermitian generator be $\op H_{\mathrm{eff}}$,
with
\begin{equation}
\ii\hbar\frac{\dd}{\dd t}\ket{\psi(t)}
 =\op H_{\mathrm{eff}}\ket{\psi(t)}.
\label{eq:nh-schro}
\end{equation}
Take the Hermitian adjoint of both sides. The imaginary unit changes
sign under complex conjugation, so
\begin{equation}
-\ii\hbar\frac{\dd}{\dd t}\bra{\psi(t)}
 =\bra{\psi(t)}\op H_{\mathrm{eff}}^\dagger.
\end{equation}
Rearranging these two equations gives
\begin{equation}
\frac{\dd}{\dd t}\ket\psi
 =-\frac{\ii}{\hbar}\op H_{\mathrm{eff}}\ket\psi,\qquad
\frac{\dd}{\dd t}\bra\psi
 =\frac{\ii}{\hbar}\bra\psi\op H_{\mathrm{eff}}^\dagger.
\end{equation}
Use the product rule, keeping both contributions:
\begin{align}
\frac{\dd}{\dd t}\braket\psi\psi
&=\left(\frac{\dd}{\dd t}\bra\psi\right)\ket\psi+
  \bra\psi\left(\frac{\dd}{\dd t}\ket\psi\right)\notag\\
&=\frac{\ii}{\hbar}\mel\psi{\op H_{\mathrm{eff}}^\dagger}\psi
 -\frac{\ii}{\hbar}\mel\psi{\op H_{\mathrm{eff}}}\psi\notag\\
&=\boxed{\frac{\ii}{\hbar}
 \mel\psi{\op H_{\mathrm{eff}}^\dagger-\op H_{\mathrm{eff}}}\psi}.
\label{eq:nh-norm-identity}
\end{align}
For $\op H_{\mathrm{eff}}=\op H_{\mathrm{eff}}^\dagger$,
the last line vanishes, recovering unitary norm conservation.
For a lossy effective model it will be nonpositive.

\subsection*{2) Concrete apparatus: two coupled optical cavities}
Consider two optical cavities, each supporting one mode of common
frequency $\omega_0$. A photon can hop coherently between them at a
real coupling frequency $J>0$. Mode $A$ additionally emits into
an external reservoir at an \emph{intensity} decay rate
$\kappa\geq0$. We study the single-photon sector with basis
$\ket A=\ket{1,0}$ and $\ket B=\ket{0,1}$.
We temporarily use units $\hbar=1$ so that the matrix entries
are angular frequencies.

\begin{figure}[tbp]
\centering
\begin{tikzpicture}[>=Latex,thick,font=\small]
 \filldraw[draw=blue!70!black,fill=blue!10,rounded corners=5pt,
           line width=1pt] (-3.5,-.65) rectangle (-1.25,.65);
 \filldraw[draw=teal!70!black,fill=teal!10,rounded corners=5pt,
           line width=1pt] (1.25,-.65) rectangle (3.5,.65);
 \node at (-2.375,.20) {cavity $A$};
 \node at (-2.375,-.21) {$\ket A,\ \omega_0$};
 \node at (2.375,.20) {cavity $B$};
 \node at (2.375,-.21) {$\ket B,\ \omega_0$};
 \draw[<->,very thick] (-1.15,0)--(1.15,0);
 \node[above] at (0,.08) {hopping $J$};
 \draw[->,very thick,red!70!black] (-2.375,-.76)--(-2.375,-1.86);
 \node[below,align=center] at (-2.375,-1.88)
      {outgoing light\\loss rate $\kappa$};
 \draw[densely dashed,rounded corners,gray]
      (-3.82,-.96) rectangle (3.82,.95);
 \node[gray!70!black,anchor=south east] at (3.72,.96)
      {two-mode subsystem};
 \node[anchor=west,align=left,font=\footnotesize] at (-.63,-1.62)
      {after emission:\\vacuum $\ket{0,0}$};
\end{tikzpicture}
\caption{Original schematic of the single-photon coupled-cavity
model. The external reservoir accepts an excitation from cavity
$A$; coherent hopping transfers it between the two cavities.
The non-Hermitian matrix below describes the \emph{conditional}
no-emission sector, rather than unconditional probability loss
in the full quantum system.}
\label{fig:nh-physical-system}
\end{figure}

The coherent part of the Hamiltonian is written in ket form as
\begin{equation}
\op H_0=\omega_0(\ket A\bra A+\ket B\bra B)
 +J(\ket A\bra B+\ket B\bra A).
\end{equation}
To obtain its matrix, insert $\bra A,\bra B$ on the left and
$\ket A,\ket B$ on the right. Because
$\braket A A=\braket B B=1$ and $\braket A B=0$,
\begin{equation}
\begin{array}{ll}
\mel A{\op H_0}A=\omega_0,&\mel A{\op H_0}B=J,\\
\mel B{\op H_0}A=J,&\mel B{\op H_0}B=\omega_0.
\end{array}
\end{equation}
Thus $H_0=\left(\begin{smallmatrix}
\omega_0&J\\J&\omega_0
\end{smallmatrix}\right)$ is Hermitian.

To represent the monitored emission channel, introduce
$\op L=\sqrt{\kappa}\ket{\mathrm{vac}}\bra A$
with $\ket{\mathrm{vac}}=\ket{0,0}$. The standard
no-click (no-jump) effective Hamiltonian is
\begin{equation}
\op H_{\mathrm{eff}}=\op H_0-\frac{\ii}{2}\op L^\dagger\op L.
\end{equation}
Let us perform the multiplication:
\begin{equation}
\begin{aligned}
\op L^\dagger\op L
&=\kappa\ket A\bra{\mathrm{vac}}
 \ket{\mathrm{vac}}\bra A\\
&=\kappa\ket A\bra A,
\end{aligned}
\end{equation}
since $\braket{\mathrm{vac}}{\mathrm{vac}}=1$.
Therefore the explicit non-Hermitian matrix is
\begin{equation}
\boxed{H_{\mathrm{eff}}=
\begin{pmatrix}\omega_0-\ii\kappa/2&J\\
J&\omega_0\end{pmatrix}.}
\label{eq:nh-explicit-matrix}
\end{equation}
Its $AA$ entry is complex, while the coherent coupling stays
real and symmetric. Indeed,
$H_{\mathrm{eff}}^\dagger-H_{\mathrm{eff}}
=\ii\kappa\ket A\bra A$.

\paragraph{Check the physical meaning of the imaginary entry.}
Write the unnormalized conditional state as
$\ket{\widetilde\psi(t)}=a(t)\ket A+b(t)\ket B$.
The Schr\"odinger equation gives the two \emph{component}
equations, before and after division by $\ii$:
\begin{equation}
\begin{aligned}
\ii\dot a&=(\omega_0-\ii\kappa/2)a+Jb,
&\dot a&=(-\ii\omega_0-\kappa/2)a-\ii Jb,\\
\ii\dot b&=Ja+\omega_0 b,
&\dot b&=-\ii Ja-\ii\omega_0 b.
\end{aligned}
\end{equation}
Using $\frac{\dd}{\dd t}|a|^2=a^*\dot a+\dot a^*a$,
we find
\begin{align}
\frac{\dd}{\dd t}|a|^2
 &=-\kappa|a|^2+\ii J(ab^*-a^*b),\\
\frac{\dd}{\dd t}|b|^2
 &=-\ii J(ab^*-a^*b).
\end{align}
The two hopping contributions cancel upon addition:
\begin{equation}
\boxed{\frac{\dd}{\dd t}\bigl(|a|^2+|b|^2\bigr)
=-\kappa|a|^2\leq0.}
\label{eq:nh-survival}
\end{equation}
This verifies Eq.~\eqref{eq:nh-norm-identity} directly.

For a normalized initial single-photon state,
$P_{\mathrm{no\ jump}}(t)=|a(t)|^2+|b(t)|^2$
is the probability of detecting \emph{no photon emission} up
to time $t$ in this monitored-channel model. Upon conditioning
on that event, the surviving state is normalized by dividing
$\ket{\widetilde\psi}$ by $\sqrt{P_{\mathrm{no\ jump}}}$.
An unconditional description includes the vacuum reached
after emission and obeys trace-preserving evolution.
Thus the decreasing no-jump norm represents a physically
identified branch of a larger probability-conserving theory.

\subsection*{3) Eigenvalues: every algebraic step}
Solve the right eigenproblem
$H_{\mathrm{eff}}\vect v_\pm=\lambda_\pm\vect v_\pm$.
The characteristic equation reads
\begin{align}
0&=\det(H_{\mathrm{eff}}-\lambda I)\notag\\
 &=\begin{vmatrix}
 \omega_0-\ii\kappa/2-\lambda&J\\
 J&\omega_0-\lambda
 \end{vmatrix}\notag\\
 &=(\omega_0-\ii\kappa/2-\lambda)
   (\omega_0-\lambda)-J^2.
\end{align}
Set $z=\lambda-\omega_0$ and expand:
\begin{equation}
(-z-\ii\kappa/2)(-z)-J^2
=z^2+\frac{\ii\kappa}{2}z-J^2=0.
\end{equation}
Complete the square:
\begin{equation}
\left(z+\frac{\ii\kappa}{4}\right)^2
=J^2-\frac{\kappa^2}{16}.
\end{equation}
We obtain the exact eigenvalues
\begin{equation}
\boxed{\lambda_\pm=\omega_0-\frac{\ii\kappa}{4}
 \pm\sqrt{J^2-\left(\frac{\kappa}{4}\right)^2}.}
\label{eq:nh-complex-spectrum}
\end{equation}
An eigenamplitude with eigenvalue
$\lambda=E-\ii\Gamma/2$ evolves as
$\exp(-\ii Et-\Gamma t/2)$, so its \emph{intensity}
decays as $\exp(-\Gamma t)$. The factor of two between
amplitude and intensity decay is essential.

\paragraph{Construct the right eigenvectors explicitly.}
The top component of $(H_{\mathrm{eff}}-\lambda I)\vect v=0$
is
\begin{equation}
(\omega_0-\ii\kappa/2-\lambda)v_A+Jv_B=0.
\end{equation}
For $J>0$ choose $v_A=1$ and solve for $v_B$:
\begin{equation}
\boxed{\vect v_\pm=
\begin{pmatrix}
1\\(\lambda_\pm-\omega_0+\ii\kappa/2)/J
\end{pmatrix}.}
\label{eq:nh-right-vectors}
\end{equation}
The second component then vanishes precisely because the
secular equation holds. These vectors may be normalized by
their Euclidean norms, but eigenvectors of distinct complex
eigenvalues need not be orthogonal. For a general non-Hermitian
matrix, left eigenvectors form a different set, and
biorthogonal expansions replace the Hermitian orthonormal
eigenbasis.

\subsection*{4) Exceptional points are more than repeated eigenvalues}
Define
\begin{equation}
\Omega=\sqrt{J^2-\left(\frac{\kappa}{4}\right)^2}.
\end{equation}
For $0\leq\kappa<4J$, the square root is real: the
frequencies $\operatorname{Re}\lambda_\pm$ split, while
$\operatorname{Im}\lambda_\pm=-\kappa/4$ are identical.
For $\kappa>4J$, the square root is imaginary: the two
frequencies coincide at $\omega_0$, while the decay rates
separate. At the threshold,
\begin{equation}
\kappa_{\mathrm{EP}}=4J,\qquad
\lambda_{\mathrm{EP}}=\omega_0-\ii J.
\label{eq:nh-ep-condition}
\end{equation}
Insert these values back into the original matrix rather than
inferring the eigenspace from the roots alone:
\begin{equation}
H_{\mathrm{eff}}-\lambda_{\mathrm{EP}}I
=J\begin{pmatrix}-\ii&1\\1&\ii\end{pmatrix}.
\end{equation}
An eigenvector satisfies $-\ii v_A+v_B=0$, hence
$v_B=\ii v_A$. With $v_A=1$,
\begin{equation}
\vect v_{\mathrm{EP}}=\begin{pmatrix}1\\\ii\end{pmatrix},\qquad
\begin{pmatrix}-\ii&1\\1&\ii\end{pmatrix}
\begin{pmatrix}1\\\ii\end{pmatrix}
=\begin{pmatrix}-\ii+\ii\\1+\ii^2\end{pmatrix}
=\begin{pmatrix}0\\0\end{pmatrix}.
\end{equation}
The second row equals $\ii$ times the first, and the matrix
has rank one. Consequently its nullspace has dimension one:
a double eigenvalue possesses only one independent
right eigenvector. This is an \emph{exceptional point} (EP).
For a Hermitian matrix, a twofold degeneracy instead has
two independent eigenvectors; eigenvalue coincidence by
itself is insufficient to identify an EP.

\begin{figure}[tbp]
\centering
\begin{tikzpicture}[x=.91cm,y=1.08cm,>=Latex,font=\footnotesize]
 \begin{scope}
 \draw[->] (0,-1.24)--(0,1.30);
 \draw[->] (0,0)--(8.32,0) node[right] {$\kappa/J$};
 \node[above right,align=left] at (.15,1.04)
      {$\operatorname{Re}(\lambda_\pm-\omega_0)/J$};
 \foreach \x in {0,2,4,6,8}{
   \draw (\x,-.035)--(\x,.035);
   \node[below] at (\x,-.08) {\x};
 }
 \draw[blue!75!black,very thick] plot[domain=0:4,samples=80]
      (\x,{sqrt(max(0,1-(\x/4)^2))});
 \draw[blue!75!black,very thick] plot[domain=0:4,samples=80]
      (\x,{-sqrt(max(0,1-(\x/4)^2))});
 \draw[blue!75!black,very thick] (4,0)--(8,0);
 \draw[densely dashed,gray] (4,-1.13)--(4,1.12);
 \node[above] at (4,1.16) {EP};
 \end{scope}
 \begin{scope}[shift={(0,-1.93)}]
 \draw[->] (0,-3.92)--(0,.34);
 \draw[->] (0,0)--(8.32,0) node[right] {$\kappa/J$};
 \node[above right] at (.15,.10)
      {$\operatorname{Im}\lambda_\pm/J$};
 \foreach \x in {0,2,4,6,8}{
   \draw (\x,-.035)--(\x,.035);
   \node[above] at (\x,.03) {\x};
 }
 \foreach \y in {-1,-2,-3}{
   \draw (-.035,\y)--(.035,\y);
   \node[left] at (-.09,\y) {\y};
 }
 \draw[orange!80!black,very thick] plot[domain=0:4,samples=80]
      (\x,{-\x/4});
 \draw[orange!80!black,very thick] plot[domain=4:8,samples=80]
      (\x,{-\x/4+sqrt(max(0,(\x/4)^2-1))});
 \draw[orange!80!black,very thick] plot[domain=4:8,samples=80]
      (\x,{-\x/4-sqrt(max(0,(\x/4)^2-1))});
 \draw[densely dashed,gray] (4,-3.83)--(4,.15);
 \end{scope}
\end{tikzpicture}
\caption{Analytic real (upper) and imaginary (lower) parts
of the passive coupled-cavity spectrum in
Eq.~\eqref{eq:nh-complex-spectrum}, with $J>0$.
The frequency branches merge at $\kappa/J=4$, while
the two attenuation-related imaginary parts separate
beyond that value. Both right eigenvectors also coalesce
at the marked EP.}
\label{fig:nh-eigenvalue-curves}
\end{figure}

\paragraph{Connection to PT symmetry, with an important distinction.}
Subtract a multiple of the identity:
\begin{equation}
H_{\mathrm{eff}}=
\left(\omega_0-\frac{\ii\kappa}{4}\right)I+
K,\qquad
K=\begin{pmatrix}-\ii\kappa/4&J\\J&\ii\kappa/4\end{pmatrix}.
\label{eq:nh-traceless}
\end{equation}
For $\mathcal P=\sigma_x$ (cavity exchange) and $\mathcal T$
(complex conjugation), direct multiplication gives
$\sigma_x K^*\sigma_x=K$. Thus the shifted matrix $K$
has PT symmetry, and its eigenvalues are $\pm\Omega$.
A physically \emph{pumped} balanced gain--loss dimer would
instead be
\begin{equation}
H_{\mathcal{PT}}=
\begin{pmatrix}\omega_0+\ii\gamma&J\\
J&\omega_0-\ii\gamma\end{pmatrix},
\qquad E_\pm=\omega_0\pm\sqrt{J^2-\gamma^2},
\end{equation}
with an EP at $|\gamma|=J$.
The apparatus of Fig.~\ref{fig:nh-physical-system} has
\emph{only a lossy cavity}; its apparent gain in $K$ is
relative to the factored-out common decay. This is an
algebraic correspondence, not a claim that cavity $B$
contains an active gain medium.

\subsection*{5) Solve the dynamics even at the exceptional point}
The eigenvector expansion breaks down at the EP, so we use
the matrix exponential directly. From Eq.~\eqref{eq:nh-traceless},
\begin{equation}
\begin{aligned}
K^2&=\begin{pmatrix}-\ii\kappa/4&J\\J&\ii\kappa/4\end{pmatrix}
     \begin{pmatrix}-\ii\kappa/4&J\\J&\ii\kappa/4\end{pmatrix}\\
&=\begin{pmatrix}J^2-\kappa^2/16&0\\0&J^2-\kappa^2/16\end{pmatrix}
=\Omega^2 I.
\end{aligned}
\end{equation}
Hence $K^{2m}=\Omega^{2m}I$ and
$K^{2m+1}=\Omega^{2m}K$. Splitting the exponential series,
\begin{align}
e^{-\ii Kt}
&=\sum_{m=0}^\infty\frac{(-1)^m(\Omega t)^{2m}}{(2m)!}I
 -\ii\sum_{m=0}^\infty
     \frac{(-1)^m(\Omega t)^{2m+1}}{(2m+1)!}
     \frac{K}{\Omega}\notag\\
&=\boxed{I\cos(\Omega t)-\ii\frac{K}{\Omega}\sin(\Omega t)}.
\label{eq:nh-exp-k}
\end{align}
Since $I$ commutes with $K$, the common diagonal shift
in Eq.~\eqref{eq:nh-traceless} factors out:
\begin{equation}
e^{-\ii H_{\mathrm{eff}}t}
=e^{-\ii\omega_0t-\kappa t/4}e^{-\ii Kt}.
\end{equation}
Prepare a single photon in cavity $A$,
$(a(0),b(0))^{\mathsf T}=(1,0)^{\mathsf T}$.
Multiplication by the first column of
Eq.~\eqref{eq:nh-exp-k} gives
\begin{align}
a(t)&=e^{-\ii\omega_0t-\kappa t/4}
 \left[\cos(\Omega t)-\frac{\kappa}{4\Omega}\sin(\Omega t)\right],
\label{eq:nh-amplitude-a}\\
b(t)&=-\ii e^{-\ii\omega_0t-\kappa t/4}
 \frac{J}{\Omega}\sin(\Omega t).
\label{eq:nh-amplitude-b}
\end{align}
At $\kappa>4J$, $\Omega$ is imaginary.
Use $\sin(\ii x)=\ii\sinh x$ and
$\cos(\ii x)=\cosh x$ to obtain real-valued
decay envelopes; the formulas themselves remain valid
by analytic continuation.

Exactly at the EP, $\Omega=0$.
The limit $\lim_{\Omega\to0}\sin(\Omega t)/\Omega=t$
and $\cos 0=1$ yield
\begin{equation}
a_{\mathrm{EP}}(t)=e^{-\ii\omega_0t-Jt}(1-Jt),
\qquad
b_{\mathrm{EP}}(t)=-\ii e^{-\ii\omega_0t-Jt}Jt.
\label{eq:nh-ep-time}
\end{equation}
The polynomial factor multiplying the exponential is a
signature of the non-diagonalizable Jordan structure.

\begin{figure}[tbp]
\centering
\begin{tikzpicture}[x=1.24cm,y=2.7cm,>=Latex,font=\footnotesize]
 \draw[->] (0,0)--(4.4,0) node[right] {$Jt$};
 \draw[->] (0,0)--(0,1.16) node[above]
      {$P_{\mathrm{no\ jump}}$};
 \foreach \x in {0,1,2,3,4}{
   \draw (\x,-.012)--(\x,.012);
   \node[below] at (\x,-.028) {\x};
 }
 \foreach \y in {0,.5,1}{
   \draw (-.02,\y)--(.02,\y);
   \node[left] at (-.07,\y) {\y};
 }
 \draw[gray!75!black,thick,densely dashed] (0,1)--(4,1);
 \draw[blue!75!black,very thick] plot[domain=0:4,samples=130]
  (\x,{exp(-\x)*((cos(deg(sqrt(.75)*\x))-
  .5/sqrt(.75)*sin(deg(sqrt(.75)*\x)))^2+
  (sin(deg(sqrt(.75)*\x))/sqrt(.75))^2)});
 \draw[orange!80!black,very thick] plot[domain=0:4,samples=130]
  (\x,{exp(-2*\x)*((1-\x)^2+\x^2)});
 \draw[gray!75!black,thick,densely dashed] (2.10,1.03)--(2.45,1.03);
 \node[anchor=west] at (2.51,1.03) {$\kappa=0$};
 \draw[blue!75!black,very thick] (2.10,.88)--(2.45,.88);
 \node[anchor=west] at (2.51,.88) {$\kappa=2J$};
 \draw[orange!80!black,very thick] (2.10,.73)--(2.45,.73);
 \node[anchor=west] at (2.51,.73) {$\kappa=4J$ (EP)};
\end{tikzpicture}
\caption{Exact no-emission survival probability
$|a(t)|^2+|b(t)|^2$ for a photon initially in cavity $A$,
computed from Eqs.~\eqref{eq:nh-amplitude-a} and
\eqref{eq:nh-amplitude-b}. The Hermitian limit conserves
norm. For $\kappa=2J$ the survival decays smoothly;
at the EP $\kappa=4J$,
$P_{\mathrm{no\ jump}}=e^{-2Jt}[(1-Jt)^2+(Jt)^2]$.
The curves use dimensionless time $Jt$.}
\label{fig:nh-survival-curves}
\end{figure}

\subsection*{6) Essential consistency checks}
\begin{enumerate}
\item \textbf{Hermitian limit.}
For $\kappa=0$ we have $\Omega=J$,
$a(t)=e^{-\ii\omega_0t}\cos Jt$,
$b(t)=-\ii e^{-\ii\omega_0t}\sin Jt$.
The norm is exactly $\cos^2 Jt+\sin^2 Jt=1$.
\item \textbf{Uncoupled lossy cavity.}
Set $J=0$ in the \emph{original} component equations.
Then $a(t)=a(0)e^{-\ii\omega_0t-\kappa t/2}$
and $b(t)=b(0)e^{-\ii\omega_0t}$.
An initial photon in $A$ has survival probability
$e^{-\kappa t}$: this checks that $\kappa$ means
\emph{intensity} decay rate.
\item \textbf{Probability in the larger system.}
The norm lost by the no-jump branch is the accumulated
probability of emission, so
$P_{\mathrm{no\ jump}}(t)+P_{\mathrm{emission}}(t)=1$
for the initial one-photon state. There is no loss of
probability in the full system plus environment.
\end{enumerate}
\paragraph{Checkpoint.}
At $\kappa=4J$, show directly that the characteristic
polynomial equals $(\lambda-\omega_0+\ii J)^2$.
Compute the nullspace of
$H_{\mathrm{eff}}-(\omega_0-\ii J)I$.
Why does the EP provide only one eigenvector?
Finally take the $\Omega\to0$ limit in
Eq.~\eqref{eq:nh-exp-k} and verify
$e^{-\ii Kt}=I-\ii Kt$ using $K^2=0$.

\noindent\textbf{Independent extension.}
In the online companion, the
\href{https://github.com/KangqiaoLiu/advanced-quantum-mechanics-lecture-note/blob/main/companion/reading.html#IC-01-09}
{guided non-Hermitian reading question}
offers an additional matrix example and a chance to
check the complex spectrum with symbolic calculation.
The hand derivation and norm balance provide the primary
physical verification.
